\section{Introduction}
\label{sec:intro}

\paragraph{Motivation: precious-metal drugs are not small molecules.}
Cisplatin and its successors---carboplatin, oxaliplatin, nedaplatin---remain
the standard of care for testicular, bladder, colorectal, and several head-and-neck
cancers forty-eight years after the 1978 FDA approval of the parent compound
\cite{footnote:cisp_history}.
The wider family of metallodrugs---Pt(II/IV), Ru(II/III), Ir(III), Os, Rh---now
spans more than thirty clinical or late-preclinical candidates
\cite{footnote:metallodrug_review}.
These chemotypes are {\em not} organic small molecules in disguise.
A chemistry-aware generative model for this space must respect at
least four axes that a generic string- or graph-sequence generator
ignores by default:
(i) the metal-ligand bond is a {\em dative} bond with a defined coordination
geometry (square-planar for Pt(II), octahedral for Ru/Ir(III)), which is not
captured by a standard covalent-bond topology;
(ii) the metal centre is {\em redox-active}, so generation-time oxidation-state
and spin-state bookkeeping matter;
(iii) ligand-exchange kinetics in physiological conditions (Cl$^-$
aquation, GSH adduction, trans-effect labilisation) mean that the {\em
thermodynamic product} and the {\em kinetic product} of a candidate
molecule can be very different molecules;
(iv) most candidate ligands are bidentate or tridentate, so the
combinatorial product space explodes with the number of donor-atom
sub-positions.
Re-purposing a token-sequence generative model by appending a
post-hoc metal-coordination filter can easily
emit chemically valid but kinetically inaccessible complexes.
A complete metallodrug generative model must encode chemistry as a
first-class citizen, not a downstream filter.

\paragraph{Existing approaches and their limits.}
Three generations of computational methods have been brought to bear on
this problem, none of them fully satisfactory.
\emph{Quantitative structure-activity relationships} (2D/3D QSAR) compress
each candidate to a fingerprint or a small set of physico-chemical
descriptors and regress against a measured endpoint
\cite{footnote:qsar_review}.
These methods are cheap and interpretable but do not propose {\em
synthesisable} molecules; they require a fixed candidate set up front.
\emph{Physics-based free-energy methods} (FEP, TI, absolute-binding
FEP) give the most defensible binding-affinity estimates and have
been applied to metallodrugs in single-pocket case studies
\cite{footnote:fep_metallo}.
They are compute-bound (tens to hundreds of nanoseconds per
candidate) and do not address the {\em construction} problem at all.
\emph{Traditional fragment-based design} assembles a lead by linking or
merging docked fragments under a medicinal-chemistry rules-of-thumb
filter; performance is dominated by the quality of the docking
function and the size of the fragment library.
\emph{Deep generative ligand design} in the last decade has produced
an impressive family of token-sequence and graph-sequence
models---REINVENT \cite{footnote:reinvent4}, GraphAF
\cite{footnote:graphaf}, latent flow-matching
\cite{footnote:latent_fm}, Junction-Tree VAEs \cite{footnote:jtvae}, and
MolDQN \cite{footnote:moldqn}---together with their
geometry-conditioned extensions (TargetDiff
\cite{footnote:targetdiff}, DiffSBDD \cite{footnote:diffsbdd},
Pocket2Mol \cite{footnote:pocket2mol}, MolDiff \cite{footnote:moldiff},
DecompDiff \cite{footnote:decompdiff}, FLOWr
\cite{footnote:flowr}) that learn to generate atoms in 3D
coordinates when a binding pocket is supplied.
Their evaluation suites, when a pocket is given, are now almost
entirely geometric: coordinate RMSD, Vina docking score, clash rate.
None of these models is formal about {\em what a molecule
is}.
They represent the molecule as the model-output (a token string, a
latent vector, an atom cloud, a graph) rather than as the object of
a calculus, and they encode synthesis routes (when they do at all)
as post-hoc reverse-synthesis filters such as Ertl-SA or
AiZynthFinder.
The combinatorial explosion of reaction space (tens of thousands of
named reactions, of which any candidate ``synthesisable molecule''
touches only a handful) plus the need for chemistry-aware
construction-time constraints therefore motivates a {\em different}
representation.

\paragraph{The gap: no framework treats molecular construction as a typed-calculus problem.}
To our knowledge, no published framework treats the generation of a
molecule as a {\em typed-calculus} problem in which the molecule is a
term in a typed $\lambda$-calculus, reaction operators are reductions,
and chemistry constraints are type predicates.
This representation gives four properties that token-sequence
generative models do not give:
(i) {\em scope} and {\em $\alpha$-equivalence} provide a
principled definition of ``the same molecule under different
variable names'' that subsumes canonical SMILES while preserving
substructural identity;
(ii) {\em $\beta$-normal form} ($\beta$-NF) gives a unique, confluent
evaluation target for ``the molecule that this construction sequence
reduces to'';
(iii) {\em typed variables} carry chemistry-aware type predicates
(aromaticity, hybridisation, ring-membership, donor-atom class) at
generation time, so synthesis-aware pruning is a typing judgement,
not a post-hoc filter;
(iv) {\em capture-avoiding substitution} guarantees that no rewrite
rule can accidentally bind a free variable of an outer context, which
is the formal analogue of ``do not break an outer ring when you
close a click reaction''.
A generation process that operates directly on terms-in-$\beta$-NF
also exposes the underlying combinatorics to a Monte-Carlo tree
search (MCTS) without requiring a differentiable surrogate.

\paragraph{Our contribution.}
We present a generator that we call the \textbf{Molecular Lambda
Calculus (MLC)}, together with a complete metallodrug-design
pipeline that uses MLC as a first-class generator.
Concretely, MLC is a \textbf{de novo typed-term MCTS, pocket-optional}
generator: the typed term is the primary object of construction and
Monte-Carlo tree search operates over $\beta$-NF terms directly;
binding-pocket conditioning enters as one of several optional oracle
channels (Vina, QuickVina, REINVENT4, PoseBusters, ADMET) rather than
as the input that defines the molecule.
\begin{itemize}
  \item \textbf{MLC as a typed calculus.} Atoms are primitive
  combinators; bonds are $\beta$-reductions; molecules are closed
  $\lambda$-terms in $\beta$-normal form. We instantiate a 9-layer
  formalism (atoms, bonds, scopes, reactions, $\alpha$-equivalence,
  $\beta$-NF, type predicates, geometry priors, oracle channels) with
  a full operational semantics (Appendix~\ref{app:betanf}). Typed
  variables carry chemistry predicates; capture-avoiding substitution
  preserves outer scopes; $\beta$-NF is unique modulo
  $\alpha$-equivalence.
  \item \textbf{Productive chemistry as a 5-click rule set.}
  The product of MLC is restricted to a curated 5-rule click-chemistry
  set (CuAAC, SPAAC, SPC, Diels--Alder, Thiol--Ene) plus a
  fallback oxidative-addition / reductive-elimination route for
  dative-bond construction. This keeps the productive space
  chemistry-realistic while remaining small enough to enumerate,
  and admits a closure theorem (Appendix~\ref{app:closure_theorem})
  that bounds reachable $\beta$-NF terms.
  \item \textbf{MetalGeometryPrior as a first-class constraint.}
  Pt(II) square-planar (coordination number 4) and Ru/Ir(III)
  octahedral (coordination number 6) coordination geometries are
  enforced as typed-predicates over the geometric term, not as
  post-hoc geometry checks.
  \item \textbf{MCTS over $\beta$-NF space.} A PUCT-guided Monte-Carlo
  tree search with virtual loss and a transposition table operates
  directly over $\beta$-NF terms, using typed-variable hits and
  click-rule applications as scoring signals.
\end{itemize}

\paragraph{Empirical evidence we report in this paper.}
All numerical claims in this paragraph are \emph{measured} in the
Mol-Metal development history; each citation points to the existing
report from which the number is taken.
\begin{itemize}
  \item \emph{Pure-Lambda validity and synthesizability on real
  pockets.} Across 5 pockets $\times$ 3 seeds $\times$ 2 arms (with
  and without a metal seed), a pure-Lambda baseline (no docking,
  no ADMET, no PoseBusters) attains
  $\mathit{validity}=1.000$, $\mathit{synthesizability}=1.000$, and
  $\mathit{metal\_compliance}=1.000$ on the metal-seed arm
  (15/15 cells) and $\mathit{validity}=1.000$,
  $\mathit{synthesizability}=1.000$, $\mathit{metal\_compliance}=0.000$
  on the organic arm (15/15 cells) \cite{footnote:wf_lambda1c}.
  \item \emph{The 5-click rule set carries search diversity.}
  An ablation that restricts the MCTS to a single click rule (CuAAC)
  produces 20\% of the search-space diversity of the full 5-rule
  set; the remaining 80\% of reachable $\beta$-NF terms require
  at least one of SPAAC, SPC, Diels--Alder, or Thiol--Ene
  \cite{footnote:wf_lambda1}.
  \item \emph{Homotype diversity is orthogonal to Tanimoto.}
  On a 10-molecule test set (5 constitutional isomers of C$_6$H$_{12}$
  + 5 unrelated drug-like molecules), the enriched homotype-distance
  metric has mean $0.1073$ on the constitutional-isomer subset
  whereas Tanimoto (Morgan r=2, 2048 bits) has mean $0.9276$; on the
  disjoint-symbol unrelated subset (cisplatin vs.\ benzene /
  naphthalene / aspirin / caffeine) the homotype distance recovers
  the disjoint-alphabet signal (cisplatin-vs-benzene homotype
  $0.5000$)
  ---so the two metrics disagree on {\em which} chemistry is
  distant rather than on the magnitude of the distance
  \cite{footnote:wf_lambda2e}.
  \item \emph{Vina vs QuickVina parity on real pockets.}
  On the 1h36 pocket, N=50 paired (same box, same exhaustiveness,
  same seed), AutoDock Vina 1.2.7 and QuickVina 2 (binary-identity
  verified upstream) agree with Pearson $r=0.9983$, Spearman
  $\rho=0.9984$, mean paired difference $+0.009\pm0.018$~kcal/mol
  ($n=46$ paired, 4 failures both attributable to a 2011-QVina
  CG0-atom-type limitation)
  \cite{footnote:round11_parity}.
  \item \emph{REINVENT4 multiproperty scoring is functional on
  gfx1101.} A real subprocess invocation of the REINVENT4 learned
  scoring CLI inside our reward aggregator returns a $[0,1]$ score
  per SMILES with mean absolute delta $0.07$ against the RDKit proxy
  on a 10-SMILES drug-like batch \cite{footnote:wf_extra2}.
  \item \emph{Pt(II) prior end-to-end pipeline is wired but not
  validated at scientific-budget.} The Round-10 micro-bench on 1h36
  confirms the harness is runnable (42.60~s full sweep, both
  prior-OFF and prior-ON settings complete, gate tests 12/12 green);
  the measured prior-effect on Vina is $\Delta = +0.000$~kcal/mol at
  $n_{\text{docked}}=4$---a statistically-inconclusive null result
  that we report honestly rather than over-projecting
  \cite{footnote:round10_e2e}.
\end{itemize}

\paragraph{Paper structure.}
The remainder of the paper is organised in seven sections.
\textbf{Section~\ref{sec:related}} surveys related work, including
de novo metallodrug design, calculus-based molecule representations,
and the geometric SBDD lineage that MLC does \emph{not} belong to,
and surfaces seven protocol-mismatch flags where the SOTA comparison
methodology is incompatible with metallodrug evaluation.
\textbf{Section~\ref{sec:method}} (a separate workflow) develops the
MLC formalism, the 9-layer architecture, the click-rule registry,
the MetalGeometryPrior, the MCTS-over-$\beta$-NF driver, and the
reward aggregator channels (Vina / QuickVina / ADMET / REINVENT4 /
PoseBusters).
\textbf{Section~\ref{sec:evaluation}} (a separate workflow) reports
the 100-pocket $\times$ 3-seed Round-13 scientific-budget run on
CrossDocked2020 (used as the cross-dock \emph{evaluation} set for
SA / QED / Novelty columns on 5 pockets, not as a training set), the
homotype-vs-Tanimoto diversity panel, the Vina
vs QuickVina parity audit, the Pt(II) prior end-to-end re-run on a
real Pt pocket, and the anticancer metric suite (TODO-15).
\textbf{Section~\ref{sec:ablation}} (a separate workflow) ablates
the 9-layer formalism, the click-rule set, the MetalGeometryPrior
weights, the MCTS virtual-loss / transposition-table components, and
the GPU kernel substitution.
\textbf{Section~\ref{sec:limitations}} (this workflow) catalogues
the architectural and empirical limits of the present system in a
checklist of seven protocol-mismatch items, six open
user-gated decisions, and the explicitly-named empirical
non-claims (Pt prior $\Delta = 0$ at this $n$; CFG decoder $0/96$ at
the unconstrained architecture; the $n_{\text{distinct}}{=}1$
singleton-attractor collapse observed when the metal seed
over-constrains the 5-click productive space; full 100-pocket
sweep not yet executed).
\textbf{Section~\ref{sec:future}} (this workflow) outlines the
Lambda $\times$ model-side coupling (deferred per
\cite{footnote:lambda_coupling}), the closure-theorem proof for the
5-click productive space, the Round-12 / Round-13 / Round-14
timeline revisions, and the open engineering questions about
fused-silu-MLP on gfx1101 wave64.
The appendix contains the full $\beta$-NF operational semantics,
the click-rule registry table, the per-pocket protocol sheets, and
the SHA-256-pinned input/output artefacts for every measured
number in the body.

\paragraph{Style and reproducibility notes.}
The paper follows an IMRaD structure with a strong introduction,
honest limitations, and an explicit future-work section.
Where a numerical claim is \emph{measured} in Mol-Metal's history,
we cite the corresponding internal report at the foot of the page
(per Digital Discovery's supplementary-URL convention); where a
claim is \emph{projected} but not yet measured, we label it
projected.
The paper, its code, and the per-cell inputs and outputs are
SHA-256-pinned; the supplementary archive is sufficient to
reproduce any individual cell without re-running the full sweep.
All round-9, round-10, round-11, and round-12 artefacts referenced
in this section are archived at
\texttt{molmetal/reports/}; the per-cell protocol sheets are at
\texttt{molmetal/reports/ultracode\_audit/}.
